Taint analyzers report path-traversal (CWE-22) warnings that are often false alarms. Using an LLM to dismiss them is attractive, but a wrongly dismissed warning is a missed vulnerability. We compare five configurations on identical warnings from Semgrep Community Edition:
\begin{itemize}
\item the symbolic baseline;
\item deterministic refinement by our own intraprocedural value analysis;
\item ordinary LLM triage;
\item LLM triage over structured program evidence with explicit obligations;
\item the same LLM responses filtered by a gate that checks the model's cited evidence against extracted facts.
\end{itemize}
The benchmark subjects are 134 held-out OWASP Benchmark for Python cases and 24 held-out diagnostic program pairs that we constructed, evaluated with \texttt{qwen3-coder-next} and \texttt{qwen3.5:9b} under a protocol frozen before testing. A case study on six real-world fix commits uses the primary model only.

In the frozen runs, deterministic refinement removed 36\% (OWASP) and 68\% (pairs) of false warnings and kept every benchmark-labelled vulnerability. Ordinary LLM triage removed up to 79\% but discarded up to 35\% of real vulnerabilities. On the pairs, structured evidence reduced these losses (2 instead of 7 with the primary model, 0 instead of 8 with the smaller one). The gate never rejected a wrong suppression and rejected 1 and 5 correct ones on OWASP. It checks that the cited guard is consistent with the dependency structure, not that the guard confines the path. An independent review found programs on which the frozen extractor wrongly proved a warning safe. We corrected the extractor, and every held-out evidence record, and hence every reported result, is unchanged. In the case study, four targets are within the task; the baseline reported the documented defect in one, a fix-related flow in another, and missed two.
